% Maintainer-only smoke test. This file is not included in user releases.
\documentclass[tocdepth3]{uestcthesis}

\begin{document}

\uestcfrontmatter

\zhabstract
中文摘要回归测试。
% 此处故意不留空行，用于验证关键词命令会自行另起一段。
\keywords{模板；测试}

\enabstract
English abstract regression test.
% Deliberately omit the blank line to verify that the command starts a paragraph.
\enkeywords{Template; Test}

\uestctableofcontents
\uestclistoffigures
\uestclistoftables
\uestcsymbols
\begingroup
\renewcommand{\arraystretch}{1.3}
\setlength{\tabcolsep}{9pt}
\noindent
\begin{tabularx}{\linewidth}{l>{\raggedright\arraybackslash}Xc}
    \textbf{符号} & \textbf{说明} & \textbf{页码} \\
    $a$ & 用于检查主要符号表的示例符号 & 1 \\
\end{tabularx}
\endgroup
\uestcacronyms
\begingroup
\renewcommand{\arraystretch}{1.3}
\setlength{\tabcolsep}{9pt}
\noindent
\begin{tabularx}{\linewidth}{%
    l
    >{\hsize=1.2\hsize\linewidth=\hsize\raggedright\arraybackslash}X
    >{\hsize=0.8\hsize\linewidth=\hsize\raggedright\arraybackslash}X}
    \textbf{英文缩写} & \textbf{英文全称} & \textbf{中文全称} \\
    UESTC & University of Electronic Science and Technology of China & 电子科技大学 \\
\end{tabularx}
\endgroup

\uestcmainmatter

\chapter{排版回归测试}
\section{标题与正文}
\subsection{三级标题}
\subsubsection{四级标题}

正文、脚注\footnote{脚注回归测试。}、文献引用\cite{lamport1994latex}与公式：
\begin{equation}
    a^2+b^2=c^2.
\end{equation}

\begin{definition}定义环境回归测试。\end{definition}
\begin{theorem}定理环境回归测试。\end{theorem}
\begin{lemma}引理环境回归测试。\end{lemma}
\begin{proposition}命题环境回归测试。\end{proposition}
\begin{corollary}推论环境回归测试。\end{corollary}
\begin{assumption}假设环境回归测试。\end{assumption}
\begin{remark}注环境回归测试。\end{remark}
\begin{proof}证明环境及证毕符号回归测试。\end{proof}

\begin{algorithm}[htbp]
    \caption{算法环境回归测试}
    \label{alg:chinese-smoke}
    \begin{algorithmic}[1]
        \Require 正整数$a$和$b$
        \Ensure $\gcd(a,b)$
        \State \Return $a$
    \end{algorithmic}
\end{algorithm}
\autoref{alg:chinese-smoke} 用于检查算法名称和编号。

\begin{figure}[htbp]
    \centering
    \rule{8cm}{3cm}
    \caption{图片标题回归测试：用于检查较长标题在限定宽度内换行后的居中、行距和末行位置}
    \fignote{注：这是一段用于版式回归测试的较长图片附注。附注需要在正文宽度内自然换行，除最后一行外各行应当保持两端对齐，最后一行自然结束；同时还需要检查附注与图片标题、后续正文之间的垂直距离是否稳定。}
\end{figure}

图片后的正文用于检查浮动体与正文间距。

% 将两个长附注分开观察，避免测试页因同时容纳两个浮动体而人为过满。
\clearpage

\begin{table}[htbp]
    \centering
    \caption{表格标题回归测试：用于检查较长标题在限定宽度内换行后的居中、行距和末行位置}
    \begin{tabular}{lc}
        \toprule
        项目 & 数值 \\
        \midrule
        示例 & 1 \\
        \bottomrule
    \end{tabular}
    \tabnote{注：这是一段用于版式回归测试的较长表格附注。附注需要在正文宽度内自然换行，除最后一行外各行应当保持两端对齐，最后一行自然结束；同时还需要检查附注与表格主体、后续章节之间的垂直距离是否稳定。}
\end{table}

\uestcacknowledgements
致谢回归测试。

\uestcprintbibliography

\uestcappendix
\chapter{附录回归测试}
附录正文。
\begin{theorem}附录定理编号回归测试。\end{theorem}
\begin{algorithm}[htbp]
    \caption{附录算法编号回归测试}
    \begin{algorithmic}
        \State \Return $0$
    \end{algorithmic}
\end{algorithm}

\uestcachievements
成果列表回归测试。

\end{document}
